\chapter{INTRODUCCIÓN}

\section{Motivación del trabajo}

La seguridad de las redes blockchain públicas descansa sobre primitivas criptográficas cuya solidez se asume computacionalmente intratable para las computadoras clásicas. En particular, Ethereum ---la segunda red blockchain por capitalización de mercado, con un volumen diario de transacciones que supera los cien mil millones de dólares--- emplea el esquema de firma digital ECDSA (\emph{Elliptic Curve Digital Signature Algorithm}) sobre la curva secp256k1 para autenticar cada transacción enviada a la red \parencite{wood2014ethereum}.

La seguridad de ECDSA se fundamenta en la dureza computacional del Problema del Logaritmo Discreto en Curvas Elípticas (ECDLP): dado un punto $Q = k \cdot G$ en la curva, donde $G$ es el generador y $k$ la clave privada, no existe un algoritmo clásico eficiente para recuperar $k$. Los mejores ataques clásicos (Pollard's rho, baby-step giant-step) requieren $O(\sqrt{n})$ operaciones, donde $n$ es el orden del grupo, lo que para secp256k1 ($n \approx 2^{256}$) implica aproximadamente $2^{128}$ operaciones: una barrera computacional infranqueable con la tecnología actual.

Sin embargo, en 1994 Peter Shor demostró que un computador cuántico de escala suficiente puede resolver el ECDLP en tiempo polinómico $O(n^3)$ \parencite{shor1994algorithms}. Esto significa que la aparición de un computador cuántico criptográficamente relevante (CRQC, \emph{Cryptographically Relevant Quantum Computer}) reduciría la seguridad de secp256k1 de $2^{128}$ operaciones clásicas a un número polinómico de operaciones cuánticas, haciendo viable la falsificación de firmas digitales y, por extensión, el robo de fondos de cualquier cuenta cuya clave pública haya sido expuesta.

Las estimaciones actuales sitúan la disponibilidad de un CRQC en un horizonte de 10 a 15 años \parencite{nistpqc2024}. Si bien este plazo puede parecer lejano, dos factores exigen una acción preventiva inmediata:

\begin{enumerate}[label=\arabic*.]
  \item \textbf{Amenaza <<harvest now, decrypt later>>}: Un adversario puede registrar hoy las transacciones firmadas que se propagan por la red ---incluyendo las claves públicas que se derivan durante la verificación--- y almacenarlas hasta que disponga de un CRQC capaz de recuperar las claves privadas correspondientes. Dado que la blockchain es un registro público e inmutable, todas las transacciones históricas están disponibles para este tipo de ataque \parencite{quantumresistance2023nature}.
  
  \item \textbf{Complejidad de la migración}: Ethereum es un sistema distribuido con cientos de miles de nodos, múltiples implementaciones de clientes (geth, Nethermind, Besu, Erigon, reth), y un ecosistema de contratos inteligentes valorado en miles de millones. Una migración criptográfica requiere coordinación a escala global y años de planificación, desarrollo, pruebas y despliegue progresivo.
\end{enumerate}

Estas consideraciones motivan el presente trabajo: demostrar la viabilidad técnica de migrar la capa de transacciones de Ethereum a criptografía post-cuántica basada en retículos, proporcionando una implementación funcional y una propuesta de estandarización.


\section{Planteamiento del problema}

El esquema ECDSA/secp256k1 no es un componente aislado de Ethereum, sino que está profundamente integrado en múltiples capas del protocolo. Se identifican cinco puntos de vulnerabilidad cuántica:

\begin{enumerate}[label=\textbf{P\arabic*.}]
  \item \textbf{Firma de transacciones}: Cada transacción incluye una firma ECDSA $(v, r, s)$ que autentica al remitente. Un CRQC podría falsificar firmas para cualquier cuenta cuya clave pública sea conocida.
  
  \item \textbf{Recuperación del remitente}: Ethereum no incluye la dirección del remitente en la transacción. En su lugar, la función \texttt{ecrecover} recupera la clave pública a partir de la firma $(v, r, s)$ y el hash del mensaje. Esta propiedad de <<recuperabilidad>> es exclusiva de ECDSA y no existe en los esquemas de firma post-cuánticos.
  
  \item \textbf{Derivación de direcciones}: Las direcciones Ethereum se derivan como $\texttt{keccak256}(\mathit{pk})[12{:}32]$, donde $\mathit{pk}$ es la clave pública no comprimida de 64 bytes (sin el prefijo \texttt{0x04}). Este esquema está acoplado al formato de clave secp256k1.
  
  \item \textbf{Precompilados criptográficos}: La EVM incluye precompilados para operaciones sobre curvas elípticas (BN254: \texttt{0x06}--\texttt{0x08}; BLS12-381: \texttt{0x0b}--\texttt{0x13}) y evaluación de puntos KZG (\texttt{0x0a}). Todas estas primitivas se basan en el problema del logaritmo discreto o en emparejamientos bilineales, vulnerables al algoritmo de Shor.
  
  \item \textbf{Consenso (Proof of Stake)}: La capa de consenso de Ethereum utiliza firmas BLS12-381 para la atestación y propuesta de bloques por parte de los validadores. BLS12-381 es igualmente vulnerable a ataques cuánticos, aunque su migración es un problema separado que requiere agregación de firmas post-cuántica ---una capacidad que aún no existe de forma estandarizada.
\end{enumerate}

En el momento de redacción de este trabajo, no existe ninguna propuesta de estandarización (EIP) que defina un tipo de transacción nativo post-cuántico. Las propuestas existentes ---EIP-7932 \parencite{eip7932} (registro de algoritmos de firma secundarios) y EIP-8051 \parencite{eip8051} (precompilado ML-DSA)--- abordan la verificación a nivel de contrato inteligente pero no la autenticación a nivel de protocolo.


\section{Objetivos}

\subsection{Objetivo general}

Demostrar la viabilidad de migrar la capa de transacciones de un cliente de ejecución Ethereum a criptografía post-cuántica basada en retículos, mediante la implementación de un tipo de transacción nativo con ML-DSA-65 y su validación en un entorno multi-nodo.

\subsection{Objetivos específicos}

\begin{enumerate}[label=\textbf{OE\arabic*.}]
  \item Simular los ataques de Shor (recuperación de clave privada ECDLP) y Grover (búsqueda de preimagen hash) sobre modelos reducidos de las primitivas criptográficas de Ethereum mediante circuitos cuánticos implementados con Qiskit.
  
  \item Diseñar un formato de transacción post-cuántica (\texttt{0x50}) compatible con la especificación EIP-2718 (\emph{Typed Transaction Envelope}), incluyendo la codificación RLP, el esquema de firma y el hash de transacción.
  
  \item Implementar la derivación de direcciones post-cuánticas basada en SHAKE-256 como alternativa a Keccak-256, garantizando la separación de dominios con las direcciones clásicas.
  
  \item Implementar un \emph{opcode} EVM (\texttt{PQHASH}, \texttt{0x21}) para computación SHAKE-256 dentro de contratos inteligentes, y un precompilado (\texttt{0x0100}) para verificación ML-DSA-65.
  
  \item Diseñar e implementar un mecanismo de consenso Proof of Authority (PoA) con sellado de bloques mediante ML-DSA-65 y rotación round-robin de validadores, como alternativa al Proof of Stake que requiere agregación BLS.
  
  \item Desarrollar un \emph{wallet} post-cuántico con interfaz de línea de comandos (CLI) y terminal interactiva (TUI) que permita la generación de claves, firma de transacciones y consulta de saldos.
  
  \item Validar la solución completa mediante una suite de pruebas extremo a extremo (E2E) ejecutada sobre un clúster de tres nodos validadores, verificando identidad de cadena, consenso, comisiones, transacciones, contratos y consistencia de estado.
  
  \item Redactar una propuesta de EIP (\emph{Ethereum Improvement Proposal}) que formalice la especificación del tipo de transacción post-cuántica para su eventual presentación a la comunidad Ethereum.
\end{enumerate}


\section{Antecedentes}

\subsection*{Estandarización post-cuántica del NIST}

En 2016, el Instituto Nacional de Estándares y Tecnología de Estados Unidos (NIST) inició un proceso de selección de algoritmos criptográficos resistentes a ataques cuánticos. Tras tres rondas de evaluación pública, en agosto de 2024 se publicaron los tres primeros estándares finales \parencite{nistfips203, nistfips204, nistfips205}:

\begin{itemize}
  \item \textbf{FIPS-203 (ML-KEM)}: Mecanismo de encapsulación de claves basado en Module-LWE (anteriormente CRYSTALS-Kyber).
  \item \textbf{FIPS-204 (ML-DSA)}: Esquema de firma digital basado en Module-LWE (antes CRYSTALS-Dilithium). Es el esquema empleado en este trabajo.
  \item \textbf{FIPS-205 (SLH-DSA)}: Esquema de firma digital basado en funciones hash (antes SPHINCS+), sin asunciones sobre retículos.
\end{itemize}

El corpus bibliográfico de este trabajo incluye las especificaciones completas de los tres estándares, así como las implementaciones de referencia de los candidatos originales del NIST (Dilithium, Falcon, Kyber, SPHINCS+).

\subsection*{Propuestas existentes en Ethereum}

EIP-2718 \parencite{eip2718} estableció en 2020 el marco de transacciones tipadas (\emph{Typed Transaction Envelope}), que permite añadir nuevos formatos de transacción sin romper la compatibilidad con los existentes. Actualmente se han asignado los tipos \texttt{0x00}--\texttt{0x04}.

EIP-7932 propone un registro de algoritmos de firma secundarios con un precompilado \texttt{sigrecover}. EIP-8051 propone precompilados específicos para verificación ML-DSA. Ambas propuestas abordan la verificación a nivel de contrato inteligente, pero ninguna define un tipo de transacción nativo, que es la contribución principal de este trabajo.

\subsection*{Trabajos relacionados en blockchain post-cuántica}

Varios trabajos académicos han explorado la integración de criptografía post-cuántica en redes blockchain. \textcite{quantumsafeblockchain2018} presentan un esquema de blockchain protegida cuánticamente. \textcite{quantumresistance2023nature} analizan la resistencia cuántica en redes blockchain desde una perspectiva general. \textcite{surveycomparison2023} ofrecen una comparativa exhaustiva de blockchains post-cuánticas y cuánticas. \textcite{quantumsafefabric2024} implementan protección post-cuántica en Hyperledger Fabric. El presente trabajo se diferencia al abordar específicamente la EVM de Ethereum con una implementación nativa a nivel de protocolo.


\section{Metodología}

El desarrollo del proyecto sigue una metodología iterativa basada en \emph{sprints} semanales, con las siguientes prácticas:

\begin{itemize}
  \item \textbf{Control de versiones}: Repositorios Git con flujo de ramas por funcionalidad (\emph{feature branches}), integración mediante \emph{pull requests} con \emph{squash and merge}.
  
  \item \textbf{Integración continua}: Cada \emph{pull request} se valida con \texttt{cargo test} (tests unitarios e integración), \texttt{cargo clippy} (análisis estático) y \texttt{cargo +nightly fmt} (formato).
  
  \item \textbf{Pila tecnológica}:
  \begin{itemize}
    \item Rust 1.95 (nodo de ejecución, \emph{wallet}, librería criptográfica)
    \item Python 3.13 + Qiskit 2.3 + Qiskit Aer 0.17 (simulación cuántica)
    \item Docker + docker-compose (despliegue multi-nodo)
    \item ratatui 0.29 + crossterm 0.28 (interfaz TUI)
  \end{itemize}
  
  \item \textbf{Estructura de repositorios}:
  \begin{itemize}
    \item \texttt{PostQuantumEVM/} --- repositorio raíz (orquestación, infraestructura, E2E)
    \item \texttt{pq-reth/} --- \emph{fork} de reth con cinco \emph{crates} post-cuánticos
    \item \texttt{pq-wallet/} --- \emph{wallet} CLI, TUI y suite E2E
    \item \texttt{ml-lattice-rs/} --- submódulo con implementación de ML-KEM y ML-DSA
    \item \texttt{qiskit-api/} --- API de simulación cuántica
  \end{itemize}
\end{itemize}


\section{Organización del documento}

El resto de la memoria se estructura en los siguientes capítulos:

\begin{itemize}
  \item \textbf{Capítulo 2 -- Marco Teórico}: Presenta los fundamentos matemáticos de la criptografía de curva elíptica, la computación cuántica (algoritmos de Shor y Grover), la criptografía basada en retículos (LWE, Module-LWE, ML-DSA), la arquitectura de Ethereum y la EVM, el cliente reth, y las funciones hash extensibles SHAKE-256.
  
  \item \textbf{Capítulo 3 -- Análisis de la Situación}: Detalla las vulnerabilidades cuánticas de Ethereum, define los requisitos funcionales y no funcionales, y presenta la solución propuesta junto con las alternativas descartadas.
  
  \item \textbf{Capítulo 4 -- Desarrollo}: Describe la implementación de cada componente: simulación cuántica (Qiskit), tipo de transacción \texttt{0x50}, primitivas criptográficas, EVM post-cuántica, consenso PoA, \emph{wallet} y infraestructura de despliegue.
  
  \item \textbf{Capítulo 5 -- Resultados}: Presenta los resultados de las simulaciones cuánticas, los tests unitarios y de integración, la validación E2E en configuración single-node y multi-nodo, el análisis de rendimiento, la comparativa con el estado del arte, y la propuesta de EIP.
  
  \item \textbf{Capítulo 6 -- Conclusiones y Trabajos Futuros}: Resume los objetivos alcanzados, identifica limitaciones del trabajo y propone líneas de trabajo futuro.
\end{itemize}
